\subsection{Self-awareness and Self-regulation in AI Systems}
\label{sec:2.3.2}
It is tempting to describe a system's self-awareness the way the term invites: a virtual assistant that can “register its internal frustration” when a request proves hard to satisfy, the way a person notices their own irritation. That description assumes an inner observer watching the system's states from somewhere — the same homunculus problem chapter~\ref{sec:4} raises for executive control, relocated to self-perception.

What self-awareness is, in the one system known to have it, is more mechanical. The neuroscientist Michael Graziano's account of consciousness treats awareness as the content of a model the brain builds of its own attention: a simplified self-description the control system uses to regulate itself, mistaken from the inside for direct acquaintance with experience \autocite{graziano2013consciousness}. The same logic extends one level up. The brain builds a model of itself as a bounded agent — a body, a history, a point of view — and that model is the self: a construction, assembled and then inhabited.

The clearest proof that the self is a model rather than a given is that it can be edited in under a minute. Seat someone at a table with one hand hidden and a rubber hand placed where their real hand would be; stroke both hands in synchrony with soft brushes. Within a minute or two, most people begin to feel the rubber hand as their own — flinching when it is threatened, misjudging the location of their real hand toward it \autocite{botvinick1998rubber}. The felt boundary of the self is a continuously computed guess, revisable the moment the evidence changes, which is the same predictive machinery chapter~\ref{sec:4} describes for perception run on the question of where the self stops.

For an AI system, “self-awareness” built this way is not a little watcher checking on the system's own states. It is a model — of the system's processing, its history, its current goal — that the system builds, consults, and updates, with the same failure modes any model has: it can be wrong, it can be manipulated by the wrong correlated inputs, and it is only ever as good as the evidence it was built from. Design for a model that can be checked and corrected, not for an inner observer whose reports have to be trusted because there is nothing behind them to check. In practice that means monitoring the self-model against the system's actual behavior for drift, refining it as evidence arrives instead of locking in a first impression, and keeping it connected to the system's model of the other parties.

\runin{What the deflationary account does not reach}
The self-model described here is machinery, and its virtue is that machinery can be inspected. An assistant that models a user's frustration, and models its own processing while doing so, is exhibiting social competence in section~\ref{sec:2.3}'s sense, and nothing above bears on whether the assistant cares how the exchange goes for the user.

The deflationary move is easy to overextend. Showing that the self is a model rather than an inner observer disposes of the homunculus. It does not show that a party capable of caring can be assembled out of models, and the two conclusions are close enough in shape to be mistaken for one. Section~\ref{sec:3.4} needs a self-model of exactly this kind, for reasons that have nothing to do with welfare — refusal that generalizes has to carry its own rationale forward and compare it against what a role has drifted into. It also needs something further, which section~\ref{sec:2.3.3} argues has no demonstrated instance without affect. A checkable self-model is a necessary component of a bearer, and it is not the hard part of one.
